% TOP_PROBLEM: {"id": 10, "title": "The Hadwiger-Nelson Problem", "category_id": 2, "category": {"id": 2, "name": "combinatorics", "display_name": "Combinatorics", "description": "Counting problems, graph theory, discrete structures.", "slug": "combinatorics", "order_index": 2, "created_at": "2026-07-31T15:26:25.671Z"}, "status": "open"}
% FIRST_POSTED: 2026-09-27
% !TeX program = lualatex
% Compile twice: lualatex 10.tex
% All references and prose are embedded; no BibTeX database or external inputs.
\documentclass[11pt,a4paper]{article}
\usepackage[margin=24mm,headheight=14pt]{geometry}
\usepackage{fontspec}
\setmainfont{Latin Modern Roman}
\setsansfont{Latin Modern Sans}
\setmonofont{Latin Modern Mono}
\usepackage{amsmath,amssymb,mathtools}
\usepackage{unicode-math}
\setmathfont{Latin Modern Math}
\usepackage{microtype}
\usepackage{enumitem,xurl,needspace}
\usepackage[dvipsnames]{xcolor}
\usepackage{hyperref}
\hypersetup{unicode=true,colorlinks=true,linkcolor=MidnightBlue,urlcolor=MidnightBlue,
 pdftitle={Research Notebook: Section 10},pdfauthor={Research notebook},bookmarksnumbered=true}
\usepackage{bookmark}
\usepackage{tocloft}
\setlength{\cftsecnumwidth}{3em}
\cftsetpnumwidth{2.5em}
\cftsetrmarg{4em}
\usepackage{titlesec}
\titleformat{\section}{\Large\bfseries\raggedright}{\thesection}{0.7em}{}
\usepackage{fancyhdr}
\pagestyle{fancy}\fancyhf{}
\fancyhead[L]{\small\sffamily Open Problems | Research notebook}
\fancyhead[R]{\small 220 | 27 September 2026}
\fancyfoot[C]{\thepage}
\setlength{\parindent}{0pt}
\setlength{\parskip}{5pt plus 1pt}
\setlength{\emergencystretch}{3em}
\setcounter{tocdepth}{1}
\setcounter{secnumdepth}{1}
\setlist[itemize]{leftmargin=3em,itemsep=5pt,topsep=4pt,parsep=0pt}
\allowdisplaybreaks
\newcommand{\N}{\mathbb{N}}
\newcommand{\Z}{\mathbb{Z}}
\newcommand{\Q}{\mathbb{Q}}
\newcommand{\R}{\mathbb{R}}
\newcommand{\C}{\mathbb{C}}
\newcommand{\F}{\mathbb{F}}
\newcommand{\A}{\mathbb{A}}
\newcommand{\PP}{\mathbb P}
\newcommand{\E}{\mathbb{E}}
\newcommand{\eps}{\varepsilon}
\newcommand{\dd}{\,\mathrm d}
\newcommand{\sref}[2]{\hyperlink{source:#1:#2}{[S#2]}}
\newcommand{\eref}[2]{\hyperlink{extra:#1:#2}{[E#2]}}
% Turn the next switch off for a shorter reading copy. The quotations preserve
% every exactTarget field, including qualifications omitted from a short summary.
\newif\ifshowcatalogue
\showcataloguetrue
\newcommand{\cataloguescope}[1]{\ifshowcatalogue
 \paragraph{Exact scope in the supplied catalogue.}
 \begin{quote}\small #1\end{quote}\fi}

% Independently compilable extraction; original section text retained.
\numberwithin{equation}{section}
\begin{document}
\setcounter{section}{0}
% ================================================================
% problemId: 10
% Canonical problem ID: 10
% exactTarget SHA-256: 4103967cfc9159dffac7a00cf99e07ac71220523e70f63d8f0fe108fc7a7115c
\clearpage
\section*{\texorpdfstring{The Hadwiger-Nelson Problem}{The Hadwiger-Nelson Problem}}\label{10}
\subsection{Definitions and mathematical statement}
The unit-distance graph of the plane has vertex set $\R^2$, with an edge between $x,y$ precisely when $|x-y|=1$. Its chromatic number is
\[
 \chi(\R^2)=\min\{k\in\N:\exists c:\R^2\to\{1,\ldots,k\},
 \ |x-y|=1\Rightarrow c(x)\ne c(y)\}.
\]
Determine this integer in ordinary ZFC. The coloring is an arbitrary function; its color classes need not be measurable, open, or periodic. Colorings of only a finite set of points provide obstructions but are not by themselves colorings of the whole plane. The exact-value target is separate from merely improving a lower or upper bound.
\par\smallskip\textit{Formulation and concept references:} \sref{10}{1}.
\cataloguescope{Determine the chromatic number of the Euclidean plane: the minimum number of colours needed to colour every point so that points at distance exactly one have distinct colours, with ordinary arbitrary colourings in ZFC and no measurability requirement.}
\subsection{Short English statement}
What is the smallest number of colors that can color the entire plane without giving the same color to two points exactly one unit apart?
% BEGIN RESEARCH ATTEMPT 220
\subsection{Research attempt}

\paragraph{Current frontier.}
De Grey's finite unit-distance obstruction establishes the lower bound
five; the ordinary upper bound seven and the remaining exact-value
question are recalled in the same work. \sref{10}{1}.
The second catalogue reference concerns regular-polygon Minkowski norms,
not a resolution for Euclidean distance. \sref{10}{2}, Abstract.
Recent work proving seven colors necessary for specified map-type
colorings imposes regularity that is expressly absent here.
\eref{10}{1}, Introduction. The present attempt therefore keeps arbitrary
ZFC colorings, and pursues exact algebraic finite witnesses rather than
assuming color classes are measurable or tiles.

\paragraph{Attack route.}
A periodic tiling might improve the upper bound, but rules out most
allowable colorings at the outset. A finite graph can improve the lower
bound, but must have exact unit edges, not merely visually close ones.
The selected route combines compactness with real algebraic geometry to
locate all necessary finite obstructions over number fields. We then test
whether dense rational-coordinate searches could implement this route;
an explicit two-coloring shows why they cannot.

\paragraph{Attempt.}
\textbf{1. Algebraic coordinates lose no finite obstruction.}
Let $H$ be a finite graph realized by distinct points $p_1,\ldots,p_N$
in the Euclidean plane, with every edge at distance one. Its realization
is an existential system over $\Q$:
\[
 (x_i-x_j)^2+(y_i-y_j)^2=1\quad(ij\in E(H)),\qquad
 (x_i-x_j)^2+(y_i-y_j)^2>0\quad(i\ne j).
\]
Additional inequalities $\ne1$ may be imposed for nonedges, although they
are unnecessary for a coloring obstruction. The real algebraic numbers
form a real closed field. Quantifier elimination for real closed fields
therefore transfers this finite existential statement from $\R$ to
$\R\cap\overline\Q$. Thus every finite real unit-distance obstruction
has an exact real-algebraic realization. This use of transfer, and not a
numerical perturbation, preserves all required unit lengths.
\eref{10}{2}, real-closed-field decision theorem.

For completeness, a graph with possibly infinitely many vertices is
$k$-colorable if all its finite subgraphs are. In ZFC this follows by
compactness of the product $\{1,\ldots,k\}^{V}$: the closed cylinder
conditions forbidding equal colors on an edge have the finite
intersection property. If every finite subgraph is colorable, their
intersection is nonempty. Apply this to the plane and combine it with
the preceding transfer. We obtain
\begin{equation}\label{10:algebraic}
 \chi(\R^2)=\chi\bigl((\R\cap\overline\Q)^2\bigr).
\end{equation}
There is even some finite real number field $K$ with
$\chi(K^2)=\chi(\R^2)$. Indeed, since the chromatic number is a finite
integer, failure of a coloring with one fewer colors has a finite witness;
all the algebraic coordinates of that witness generate one finite
extension $K/\Q$. Inclusion gives the other inequality. This argument
asserts existence, not a degree bound for $K$ or a bound on the number of
vertices of the witness.

\textbf{2. An exact rational-coordinate obstruction to the search strategy.}
Despite density of $\Q^2$ in the plane, its unit-distance graph is
bipartite. Here is a proof with the denominator issue included. Put
\[
 R=\Z_{(2)}=\{a/b\in\Q:b\text{ odd}\}.
\]
Every rational unit vector $(u,v)$ belongs to $R^2$. Otherwise choose a
common denominator with a positive power of two and clear its odd part.
After removing the common power of two from the numerators, the equation
$u^2+v^2=1$ would say that the sum of two integer squares, at least one
odd, is divisible by four. It is congruent to either one or two modulo
four, a contradiction. Reducing $u^2+v^2=1$ modulo two now gives
\begin{equation}\label{10:parity}
 \overline u+\overline v=1\quad\hbox{in }R/2R\cong\F_2.
\end{equation}

Choose one representative $r_C$ for each coset $C$ of $R^2$ in $\Q^2$.
For $p\in C$, color $p$ by the sum modulo two of the coordinates of
$p-r_C$. A unit edge stays within one coset, because its displacement is
in $R^2$, and \eqref{10:parity} says that its endpoint colors differ.
This is a well-defined arbitrary two-coloring. Edges exist, so
\[
 \chi(\Q^2)=2.
\]
Consequently no amount of rational-grid refinement can find even a
three-chromatic unit-distance graph with exact rational coordinates.
Rounding an algebraic realization to rational points necessarily breaks
some crucial unit equations. Topological density is not density in the
space of exact unit-distance realizations.

\textbf{3. What an algebraic search can certify.}
One can enumerate finite graphs, decide their $k$-colorability by finite
exhaustion, and decide the existential real realization system above.
If $\chi(\R^2)>k$, compactness guarantees that this search eventually
finds a finite obstruction. Alternatively one can enumerate exact
algebraic points and test edges using real-algebraic arithmetic. Thus
lower-bound witnesses are discoverable in principle.

Neither procedure tells us when to stop searching for a nonexistent
witness. The assertion that every finite graph in this geometric class
is $k$-colorable is universal over an unbounded family, not a single
real-closed-field sentence. Quantifier elimination for \emph{each}
finite graph does not decide the chromatic number of the infinite plane.
The number-field observation in \eqref{10:algebraic} likewise has no
uniform bound on degree, height, or graph size.

Two exact filters are useful but insufficient. Distinct vertices of a
unit-distance graph have at most two common neighbors, because two
circles of radius one have at most two intersection points. Hence $K_{2,3}$
cannot occur even as a noninduced subgraph. Also, a vertex-critical
$k$-chromatic graph has minimum degree at least $k-1$: otherwise a
$(k-1)$-coloring after deleting a low-degree vertex extends to it.
A hypothetical six-chromatic obstruction can therefore be chosen with
minimum degree five and no $K_{2,3}$. These combinatorial restrictions do
not establish geometric realizability of the surviving candidates.

An attempted upper-bound transfer fails in the opposite direction:
the two-coloring of $\Q^2$ has no continuity property and cannot be
extended by taking limits along rational points. Even imposing a
choice of limiting color would not preserve all unit edges. The exact
argument proving \eqref{10:algebraic} uses algebraic transfer plus
compactness, not density or continuity.

\paragraph{Stress test.}
No measure is assigned to color classes. Compactness is used in ordinary
ZFC, matching the entry. Algebraic realizations keep vertices distinct;
forgetting that condition could collapse a purported graph obstruction.
The coloring proof for $\Q^2$ treats every denominator through cosets of
$\Z_{(2)}^2$, rather than assuming that all rational points have odd
denominators. The circle-intersection filter concerns distinct centers;
coincident centers are excluded by the realization inequalities.

Finite witnesses give lower bounds only. The present argument constructs
neither a six-chromatic graph nor a five- or six-coloring of the whole
plane. The map-type restrictions in \eref{10}{1} cannot replace either
of these missing objects.

\paragraph{Outcome.}
\textbf{Outcome: Exploratory approach with unresolved gap.}

The attempt gives an exact number-field search framework and proves that
rational-coordinate searches are structurally incapable of discovering
the relevant obstructions. It also identifies necessary critical-graph
filters. These deductions are useful diagnostics, not a new numerical
bound or a claim of priority for compactness or rational-plane coloring.

A resolution still requires an unrestricted upper coloring or an exact
finite obstruction strong enough to close the remaining integer gap.
No bounded stopping criterion for the algebraic search has been obtained.
% END RESEARCH ATTEMPT 220
\subsection{Sources}
\begin{itemize}\raggedright
\item[\textnormal{[S1]}]\hypertarget{source:10:1}{} de Grey, The chromatic number of the plane is at least 5 (2018).
\url{https://arxiv.org/abs/1804.02385}.
{\small\textit{Catalogue formulation source.}}
\item[\textnormal{[S2]}]\hypertarget{source:10:2}{} Status source for Chromatic Number of the Plane (Hadwiger--Nelson Problem).
\url{https://arxiv.org/abs/2108.12861}.
{\small\textit{Catalogue research/status reference; not a proof endorsement.}}
% BEGIN ADDED REFERENCES 220
\item[\textnormal{[E1]}]\hypertarget{extra:10:1}{}
G. Sokolov and V. Voronov, \emph{On the chromatic number of the plane for map-type
colorings}, arXiv:2502.01958, Introduction and the hypotheses of its
map-type coloring theorems. \url{https://arxiv.org/abs/2502.01958}.
Consulted 27 September 2026; no inference for arbitrary ZFC colorings is made.
\item[\textnormal{[E2]}]\hypertarget{extra:10:2}{}
A. Tarski, \emph{A Decision Method for Elementary Algebra and Geometry},
2nd ed., University of California Press, 1951, the decision and elimination
method for ordered real closed fields.
\url{https://www.rand.org/pubs/reports/R109.html}.
Used for the finite realization transfer, not for a decision procedure
for the infinite unit-distance graph.
% END ADDED REFERENCES 220
\end{itemize}

\end{document}
